\documentclass[10pt,a4paper]{amsart}
\usepackage[margin=19mm]{geometry}
\usepackage{amsmath,amssymb,mathtools}
\usepackage[scr=rsfs,cal=euler]{mathalfa}
\usepackage{xfrac}
\usepackage[unicode,hidelinks,bookmarks=false]{hyperref}
\newcommand{\ie}{i.e.\ }
\newcommand{\cf}{cf.\ }
\newcommand{\resp}{respectively\ }
\newcommand{\Subsection}{\subsection}
\providecommand{\nobreakdash}{\nobreak\hbox{-}}
\setlength{\emergencystretch}{2em}
\multlinegap0pt
\usepackage{bm}
\usepackage{bbm}
\usepackage{dsfont}
\usepackage{xspace}
\usepackage{cancel}
\usepackage{mathtools}
\usepackage{amsthm}
\makeatletter
\@ifundefined{theorem}{\newtheorem{theorem}{Theorem}}{}
\@ifundefined{defn}{\newtheorem{defn}[theorem]{Definition}}{}
\@ifundefined{lemma}{\newtheorem{lemma}[theorem]{Lemma}}{}
\@ifundefined{prop}{\newtheorem{prop}[theorem]{Proposition}}{}
\makeatother
\newcommand{\Be}{\mathcal B}
\newcommand{\Q}{{\mathbb Q}} 
\newcommand{\Z}{{\mathbb Z}} 
\newcommand{\allroots}{\Phi}
\newcommand{\al}{\alpha}
\newcommand{\be}{\beta} 
\newcommand{\coplanar}{coplanar\xspace} 
\newcommand{\macd}[2]{j^{#1}_{#2}({\rm{sgn}})} 
\newcommand{\posroots}{\Phi^+} 
\newcommand{\se}{\subseteq} 
\title[Orthogonal roots, Macdonald representations, and quasiparabo]{Orthogonal roots, Macdonald representations, and quasiparabolic sets}
\author{R.M. Green, Tianyuan Xu}
\date{Source: arXiv:2409.01948v2 (CC BY 4.0)}
\begin{document}
\maketitle

\noindent\emph{Source and licence.} Orthogonal roots, Macdonald representations, and quasiparabolic sets. Version arXiv:2409.01948v2 (CC BY 4.0), distributed under the Creative Commons Attribution 4.0 International licence, as recorded in the arXiv licence metadata for this exact version and shown on the abstract page https://arxiv.org/abs/2409.01948v2. Full rights notice: licenses/arxiv-2409.01948-CC-BY-4.0.txt.\par\smallskip

\noindent\textit{Editorial scope.} Counted unit: the Theorem whose source label is \texttt{thm:positive} in the article (source0.tex lines 2306--2326, source label \texttt{thm:positive}), delivered together with the article's own complete original proof of it (source0.tex lines 2328--2404, one \texttt{proof} environment of 77 lines). No mathematics is written, repaired or re-derived: statement and proof are the article's own text line for line. Required context retained uncounted: prop:four (lines 991--1004); def:nestcross (lines 1063--1078); thm:cna (lines 1080--1119); def:stats (lines 1482--1502); prop:miracle (lines 1580--1607); thm:indep (lines 2209--2220); lem:pm1 (lines 2293--2296). Every definition, display and label of those lines is retained unchanged. Omitted from this package and not redistributed: 1--990, 1005--1062, 1079--1079, 1120--1481, 1503--1579, 1608--2208, 2221--2292, 2297--2305, 2327--2327, 2405--4334. Nothing inside a retained excerpt is omitted or altered: every retained body is delivered exactly as the article's lines are written, so no omission receipt of any kind accompanies this package. Citations: the retained excerpts carry no citation command at all, so no bibliography entry of the article is transcribed and no reference is redistributed. Reference adaptation: none was needed and none was made; every reference inside the delivered unit resolves inside it, so no unresolved reference or citation is produced. Printed slips kept literal: no printed word, symbol, hypothesis or number of the counted statement or of its proof was altered. Layout and class: the article's own class is \texttt{amsart} with packages including babel, amsmath, graphicx, hyperref, amsthm, amssymb, amsfonts, amscd, showkeys, fullpage, subfig, enumitem, tikz, xspace; the wrapper is the lane's pinned \texttt{amsart} editorial wrapper at 10pt on A4, which restates the theorem-like environments the retained text needs and adds the article's own macro definitions that the retained text needs (\texttt{\textbackslash Be}, \texttt{\textbackslash Q}, \texttt{\textbackslash Z}, \texttt{\textbackslash al}, \texttt{\textbackslash allroots}, \texttt{\textbackslash be}, \texttt{\textbackslash coplanar}, \texttt{\textbackslash macd}, \texttt{\textbackslash posroots}, \texttt{\textbackslash se}), with extra wrapper packages (bm, bbm, dsfont, xspace, cancel, mathtools). The delivered numbering is the unit's own. Assets: no retained span contains a figure environment, a graphics-inclusion command, a PDF-page inclusion, a table or a TikZ picture; no image file is copied into this package, the package declares no asset dependency and no image file is redistributed with it. Licence: version arXiv:2409.01948v2 of this article is distributed under the Creative Commons Attribution 4.0 International licence, as recorded in the arXiv licence metadata for this exact version and shown on the abstract page https://arxiv.org/abs/2409.01948v2. A full rights notice is delivered at licenses/arxiv-2409.01948-CC-BY-4.0.txt. Characters: every retained excerpt is pure ASCII, so the delivered engine, XeLaTeX with its default Latin Modern text font, reports no missing character.
\begin{prop}\label{prop:four}
Let $W$ be a Weyl group of type $E_7$, $E_8$, or $D_n$ for $n$ even, let $\al$
be a root, and let $R$ be a maximal set of orthogonal positive roots. Suppose
that neither $\al$ nor $-\al$ is an element of $R$.
\begin{enumerate}
    \item[{\rm (i)}] The root $\al$ is orthogonal to all but precisely four
     elements $Q = \{\be_1, \be_2, \be_3, \be_4\}$ of $R$.  The elements of $Q$
     form a \coplanar quadruple, and we have $ 2\al = \pm \be_1 \pm \be_2 \pm
     \be_3 \pm \be_4$ for suitable choices of signs.
 \item[{\rm (ii)}] Let $\Psi$ be the $D_4$ subsystem associated to $Q$. Then we
     have $\alpha\in \Psi$, and the set $s_\al(Q)=\{s_\al(\be_i):1\le i\le 4\}$
     is also a \coplanar quadruple whose associated $D_4$-subsystem is $\Psi$.
\end{enumerate}
\end{prop}

\begin{defn}\label{def:nestcross}
Let $Q = \{\be_1, \be_2, \be_3, \be_4\}$ be a \coplanar quadruple of positive
orthogonal roots, let $\Psi$ be the $D_4$-subsystem associated to $Q$, let
$\leq$ be the partial order on $\Psi$ relative to the induced simple roots of
$\Psi$, and let $\gamma$ be the root $(\be_1 + \be_2 + \be_3 + \be_4)/2$. We say
that $Q$ is
\begin{enumerate}
    \item[{\rm (i)}]{a {\it crossing} if $\be_i \leq \gamma$ for all $i$ and $Q$
     contains the unique $\leq$-maximal element of $Q \cup (-s_\gamma(Q))$;}
 \item[{\rm (ii)}]{a {\it nesting} if $\be_i \leq \gamma$ for all $i$ and $Q$
     contains the unique $\leq$-minimal element of $Q \cup (-s_\gamma(Q))$;}
 \item[{\rm (iii)}]{an {\it alignment} otherwise.}
\end{enumerate} We also call each crossing, nesting, and alignment a
    \emph{feature} of \emph{type C}, \emph{type N}, and \emph{type A},
    respectively.
\end{defn}

\begin{theorem}\label{thm:cna}
Let $\allroots$ be a root system for a Weyl group of type $E_7$, $E_8$, or $D_n$
for $n$ even.  Let $Q$ be a \coplanar quadruple of positive roots of
$\allroots$, let $\Psi$ be the associated $D_4$-subsystem, and let $\Psi^+ =
\Psi \cap \posroots$ be the induced positive system of $\Psi$.
\begin{enumerate}
\item[{\rm (i)}]
The set $\Psi^+$ contains precisely three distinct quadruples of mutually
orthogonal roots. These quadruples are pairwise disjoint and partition $\Psi^+$.

\item[{\rm (ii)}]
The three quadruples of orthogonal roots in $\Psi^+$  are all coplanar. Among
them there is exactly one crossing, $\Psi^+_C$, exactly one nesting, $\Psi^+_N$,
and exactly one alignment, $\Psi^+_A$. In particular, the quadruple $Q$ cannot
be both a crossing and a nesting, and the three conditions in Definition
\ref{def:nestcross} are mutually exclusive.

\item[{\rm (iii)}]
Each quadruple in $\{\Psi^+_C,\Psi^+_N, \Psi^+_A\}$ uniquely determines both of
the other two.

\item[{\rm (iv)}]
{If $R$ is a set of mutually orthogonal roots that is disjoint from $\Psi$, then
either each of the three sets $\{R \cup \Psi^+_C, \ R \cup \Psi^+_N, \ R \cup
\Psi^+_A\}$ consists of mutually orthogonal roots, or none of them does.}

\item[{\rm (v)}]
The crossing $\Psi^+_C$ contains no root from the induced simple system of
$\Psi$.

\item[{\rm (vi)}]
For each $x\in \{C,N,A\}$, let $\gamma_x$ be the product of the roots in
$\Psi^+_x$, and let $\sigma(\gamma_x)$ be the sum of the components of $\gamma_x$. Then we have $\sigma(\gamma_A) < \sigma(\gamma_N) =
\sigma(\gamma_C)$ and $\gamma_C = \gamma_N + \gamma_A.$ Moreover, if $\al$ is
any component in one of the three $4$-roots $\gamma_C,\gamma_N$ and $\gamma_A$,
then the reflection $s_\al$ sends the other two $4$-roots to each other; for
example, if $\al\in \Psi^+_C$ then $s_\al$ sends $\gamma_N$ and $\gamma_A$ to
each other.
\end{enumerate}
\end{theorem}

\begin{defn}\label{def:stats} Let $W$ be a Weyl group of type $E_7, E_8$ or
      $D_{n}$ for $n$ even.  Let $R$ be a set of mutually orthogonal roots of
      $W$, and let $\beta$ be a positive $n$-root of $W$.
      \begin{enumerate}
          \item[(i)] We define the {\it crossing number} $C(R)$, the
               \emph{nesting number} $N(R)$, and the \emph{alignment number}
               $A(R)$ of $R$ to be the numbers of crossings, nestings, and
               alignments contained in $R^+$, respectively.
           \item[(ii)] We define the \emph{type} of $R$ to be the monomial
      $A^{A(R)}C^{C(R)}N^{N(R)}$, and define the {\it level} $\lambda(R)$ of $R$
      to be $C(R) + 2N(R)$.
  \item[(iii)] If $R$ is a maximal orthogonal set of roots, then we say that $R$
      is \emph{noncrossing}, \emph{nonnesting}, and \emph{alignment-free} if
      $C(R)=0$, $N(R)=0$ and $A(R)=0$, respectively; we also call $R$ {\it
      maximally crossing}, {\it maximally nesting}, or {\it maximally aligned}
      if we have $N(R)=A(R)=0$, $C(R)=A(R)=0$, or $C(R)=N(R)=0$, respectively.
\end{enumerate} We also apply all the above definitions to $\beta$ by applying them to
    the set of components of $\beta$. (Note that the
    definitions of type in (ii) and in Definition \ref{def:nestcross} are
consistent.)
\end{defn}

\begin{prop}\label{prop:miracle}
Let $W$ be a Weyl group of type $E_7$, $E_8$, or $D_{n}$ with $n$ even, and let
$R$ be a maximal set of  orthogonal positive roots of type $A^p C^q N^r$. Let
$\lambda=C+2N$ be the level function from Definition \ref{def:stats}.
\begin{enumerate}
    \item[{\rm (i)}]{If $\al_i$ is a simple root, then either $\al_i \in R$, or
         $\lambda(s_{\al_i}(R)) \ne \lambda(R)$.  If $\lambda(s_{\al_i}(R)) >
         \lambda(R)$, then we have $\lambda(s_{\al_i}(R)) = \lambda(R) + 1$, and
         either (1) $s_{\al_i}$ moves an $A$ to a $C$ and $s_{\al_i}(R)$ has
         type $A^{p-1}C^{q+1}N^r$, or (2) $s_{\al_i}$ moves a $C$ to an $N$ and
     $s_{\al_i}(R)$ has type $A^pC^{q-1}N^{r+1}$.}
 \item[{\rm (ii)}]If $Q$ is an alignment in $R$, $Q'$ is the corresponding
     crossing quadruple, and $R'=(R\setminus Q)\cup Q'$,  then $d = \lambda(R')
     - \lambda(R)$ is a positive odd number.  If there is a positive root $\al$
     such that $s_\alpha(Q)=Q'$ but no positive root $\al'<\al$ moves an $A$ in
     $R$ to
     a $C$, then $R'$ has type $A^{p-1}C^{q+1}N^{r}$.
 \item[{\rm (iii)}]If $Q$ is a crossing in $R$, $Q'$ is the corresponding
      nesting quadruple, and $R'=(R\setminus Q)\cup Q'$, then $A(R) = A((R
      \backslash Q) \cup Q')$ and $d = \lambda(R') - \lambda(R)$ is a positive
      odd number.  If there is a positive root $\al$ such that $s_\alpha(Q)=Q'$
      but no positive root $\al'<\al$ moves a $C$ in $R$ to an $N$, then $R'$
      has type $A^{p}C^{q-1}N^{r+1}$.
  \item[{\rm (iv)}] If $Q$ is an alignment in $R$, $Q'$ is the corresponding
      nesting quadruple, and $R'=(R\setminus Q)\cup Q'$, then $d = \lambda((R
      \backslash Q) \cup Q') - \lambda(R)$ is a strictly positive even number.
\end{enumerate}
\end{prop}

\begin{theorem}\label{thm:indep}
Let $W$ be a Weyl group of type $E_7$, $E_8$, or $D_n$ for $n$ even. Let
$\macd{\Phi}{nA_1}$ be the Macdonald representation of $W$.
\begin{enumerate}
    \item[{\rm (i)}]{The nonnesting positive $n$-roots form a $\Q$-basis for
     $\macd{\Phi}{nA_1}$.}
 \item[{\rm (ii)}]{The noncrossing positive $n$-roots form a $\Q$-basis for
     $\macd{\Phi}{nA_1}$.}
 \item[{\rm (iii)}]{The alignment-free positive $n$-roots span
     $\macd{\Phi}{nA_1}$.}
\end{enumerate}
\end{theorem}

\begin{lemma}
    \label{lem:pm1}
{If $\be$ and $\lambda \be$ are both $n$-roots for some scalar $\lambda$, then we must have $\lambda = \pm 1$.}
\end{lemma}

\begin{theorem}\label{thm:positive}
Let $W$ be a Weyl group of type $E_7$, $E_8$, or $D_n$ for $n$ even, and let
$\Be$ be the set of noncrossing positive $n$-roots.
\begin{enumerate}
    \item[{\rm (i)}]{Every $n$-root is a $\Z$-linear combination of elements of
     $\Be$, with coefficients of like sign. This sign is positive if the
 $n$-root is positive, and is negative if the $n$-root is negative.}
\item[{\rm (ii)}]{A positive $n$-root is noncrossing if and only if it is not a
     positive linear combination of other positive $n$-roots.}
 \item[{\rm (iii)}] Define $\gamma\le_\Be\gamma'$ for positive $n$-roots
     $\gamma=\sum_{\be \in \Be} c_\be \be$ and $\gamma'=\sum_{\be\in \Be}c'_\be
     \be$ whenever $c_\be \leq d_\be$ for all $\be \in \Be$. Then $\le_\Be$ is a
     partial order on the set $X$ of positive $n$-roots. The maximally crossing
     element $\theta_C$ is the unique maximal element of $X$ with respect to
     $\le_\Be$.
 \item[{\rm (iv)}]{If $\gamma \in \Be$ and $\al_i$ is a simple root, then we
     have $$ s_{\al_i}(\gamma) = \begin{cases} -\gamma & \text{\ if\ } \al_i |
         \gamma;\\ \gamma + \gamma' & \text{\ otherwise,\ for\ some\ } \gamma'
         \in \Be {\text{\ such\ that\ }} \al_i | \gamma'.\\
 \end{cases} $$}
\end{enumerate}\end{theorem}

\begin{proof}
Let $\be$ be a positive $n$-root. By the proof of Theorem \ref{thm:indep}, the
result of applying reductions of the form $\gamma \gamma_C = \gamma \gamma_N +
\gamma \gamma_A$ to $\be$ until this is no longer possible has the effect of
expressing $\be$ as a positive integer linear combination of noncrossing
$n$-roots, and this procedure will always terminate after finitely many steps.
This proves (i) for positive $n$-roots, and the statement for negative $n$-roots
follows because $n$-roots occur in positive-negative pairs.

If $\be$ is a positive $n$-root that contains a crossing, then $\be$ is a
positive linear combination of other positive $n$-roots by applying the
reduction rule in the first paragraph. Conversely, suppose that $\be$ is a
noncrossing $n$-root and that $\be = \sum_i \lambda_i \be_i$, where $\lambda_i >
0$ and $\be_i$ is a positive $n$-root that is different from $\be$  for each
$i$. Part (i) implies that each of the $\be_i$ is a positive linear combination
of noncrossing $n$-roots. Because no cancellation can occur in the sum, Theorem
\ref{thm:indep} (ii) implies that this is only possible if each $\be_i$ is a
multiple of $\be$.  Collecting terms, we then have $\be = \lambda \be_i$. Lemma
\ref{lem:pm1} and the assumption that $\be_i$ is positive then imply that
$\lambda=1$ and $\be = \be_i$, which is a contradiction.

The relation $\le_\Be$ in (iii) is clearly a partial order on $X$.  Since $X$ is
finite, it contains at least one maximal positive $n$-root with respect to
$\le$. To prove (iii), it then suffices to show that for every  $n$-root
$\gamma\in X$ not equal to $\theta_C$, there is an element $\gamma'\in X$ such
that $\gamma<_\Be\gamma'$.  Let $\gamma\in X$ be an $n$-root not equal to
$\theta_C$, so that we can factorize $\gamma$ as $\gamma_1 \gamma'$, where
$\gamma'$ is either an alignment or a nesting. In either case, Theorem
\ref{thm:cna} (vi) implies that there exists a reflection $s_\al$ such that
$s_\al(\gamma')$ is a crossing, and that we have $s_\al(\gamma') = \gamma' +
\gamma''$, where $\gamma''$ is a nesting if $\gamma'$ is an alignment, or vice
versa. We then have $$ s_\al(\gamma) = \gamma + \gamma_1\gamma'' ,$$ where
$\gamma_1\gamma''$ is also a positive $n$-root. If we write  the $n$-root
$s_\al(\gamma)$ as $s_\al(\gamma) = \sum_{\be \in \Be} e_\be \be$, then it
follows from (i) that $c_\be \leq e_\be$ for all $\be \in \Be$. It follows that
$\gamma<_\Be s_\al(\gamma)$, proving (iii).

In the situation of (iv), it is immediate that if $\al_i$ is a component of
$\gamma$, then $s_{i}(\gamma) = -\gamma$.  Suppose from now on that this is not
the case, and let $A^pN^r$ be the type of $\gamma$.  Let $R$ be the set of
components of $\gamma$, let $Q\se R$ be the coplanar quadruple moved by
$\al_{i}$, and let $\Psi$ be the $D_4$-subsystem of $Q$. Then the sets $Q,
Q'=s_i(Q)$ and $Q''=\Psi^+\setminus (Q\cup Q')$ are the three distinct coplanar
quadruples partitioning the induced positive system by Proposition
\ref{prop:four} (ii) and Theorem \ref{thm:cna}. Let $\Psi^+_A, \Psi^+_C$ and
$\Psi^+_N$ be the alignment, crossing, and nesting in $\Psi^+$, respectively.
Then since $\al_i$ is a simple root, Proposition \ref{prop:miracle} (i) implies
that we must have one of the following two situations:
\begin{enumerate}
    \item[(1)] $Q=\Psi^+_A$, $Q'=\Psi^+_C$, $s_i(\gamma)$ has type
         $A^{p-1}CN^r$, and $Q''=\Psi^+_N$;
     \item[(2)] $Q=\Psi^+_N$, $Q'=\Psi^+_C$, $s_i(\gamma)$ has type
         $A^pCN^{r-1}$, and $Q''=\Psi^+_A$.
 \end{enumerate} Let $\gamma_x=\prod_{\beta\in x}\beta$ for all $x\in
    \{Q,Q',Q''\}$, and write $\gamma=\gamma_1\gamma_Q$. Then we have
    $\gamma_{Q'}=\gamma_Q+\gamma_{Q''}$, and thus,
    \[s_i(\gamma)=\gamma_1\gamma_{Q'}=\gamma+\gamma_1\gamma_{Q''} \] in both of
    the above cases. We have $\al_i\in \Psi$ by Proposition \ref{prop:four}
    (ii), and $\al_i$ must lie in the induced simple system of $\Psi$ since it
    is simple. Theorem \ref{thm:cna} (v) then implies that $\al\notin Q'$, and
    we have $\al_i\notin Q$ by assumption, so we have $\al\in Q''$. It follows
    that $\al$ divides the $n$-root $\gamma''=\gamma_1\gamma_{Q''}$.

It remains to prove that $\gamma''$ is noncrossing.  We treat case (1) first.
If $\al$ is any root that is minimal in $Q$, then $\al$ moves the crossing $Q'$
to the nesting $Q''$ by Theorem \ref{thm:cna} (vi), so that
$\gamma''=s_\al(s_i(\gamma))$.  Since $s_i(\gamma)$ has type $A^{p-1}CN^r$, $Q'$
is the unique crossing in $s_i(\gamma)$; therefore any root that moves a $C$ in
$s_i(\gamma)$ to an $N$ must move $Q'$, and it must move $Q'$ to $Q''$ by
Proposition \ref{prop:four} (ii). Together with Theorem \ref{thm:cna} (vi), this
further implies that any root moving $Q'$ to $Q''$ must come from $Q$, so it
follows that $\al$ is minimal among roots moving a $C$ in $s_i(\gamma)$ to an
$N$. It follows from Proposition \ref{prop:miracle} (iii) that
$\gamma''=s_\al(s_i(\gamma))$ has type $A^{p-1}N^{r+1}$; therefore $\gamma''$ is
noncrossing. A similar argument shows that $\gamma''$ has type $A^{p+1}N^{r-1}$
in case (2), so $\gamma''$ is noncrossing in both cases.
\end{proof}

\end{document}
